\subsection{Differentiability and Linear Approximation}

\begin{theorem} \label{thm:bounded-derivative-convex-mean-value}
\hlt{Mean-Value Inequality on Convex Sets}\\
Let $E\subseteq\R^n$ be open and convex, and let $\f:E\to\R^m$ be differentiable.
\begin{enumerate}[label=(\roman*)]
\setlength{\itemsep}{0pt}
\item If $M\geq0$ and $(\forall\,\x\in E)\,\normop{D\f(\x)}\leq M$, then $(\forall\,\a,\b\in E)\,\norm{\f(\b)-\f(\a)}_2\leq M\norm{\b-\a}_2$.
\item If $(\forall\,\x\in E)\,D\f(\x)=0$, then $\f$ is constant on $E$.
\end{enumerate}
\begin{proof}
\begin{enumerate}[label=(\roman*)]
\setlength{\itemsep}{0pt}
\item Fix $\a,\b\in E$, $\gamma(t)\defeq(1-t)\a+t\b$, $\mathbf g\defeq\f\circ\gamma$. Then $\gamma([0,1])\subseteq E$ by convexity, and $\mathbf g'(t)=D\f(\gamma(t))(\b-\a)$ by Theorem \ref{thm:chain-rule-several-variables}.\\
Hence $\norm{\mathbf g'(t)}_2\leq M\norm{\b-\a}_2$ by the operator-norm hypothesis. Therefore $\norm{\f(\b)-\f(\a)}_2\leq M\norm{\b-\a}_2$ by Theorem \ref{thm:vector-valued-mean-value-inequality}.
\item If $(\forall\,\x\in E)\,D\f(\x)=0$, then $(\forall\,\a,\b\in E)\,\norm{\f(\b)-\f(\a)}_2=0$ by part (i) with $M=0$; $(\forall\,\a,\b\in E)\,\f(\a)=\f(\b)$.
\end{enumerate}
\hlt{Structural Commentary}: Convexity retains the segment; the chain rule transfers the bound.\\
\hlt{Key Insight}: A vector mean-value bound converts derivative control into secant control.
\end{proof}
\end{theorem}

\begin{definition} \label{def:c1-map} \hlt{$C^1$ Mapping}\\
Let $E\subseteq\R^n$ be open, and let $\f:E\to\R^m$ be differentiable.\\
The map $\f$ is \hlt{continuously differentiable}, or $C^1$, if $D\f:E\to\LL(\R^n,\R^m)$ is continuous in the operator-norm metric.
\end{definition}

\begin{perspective} \label{per:c1-operator-norm-continuity}
\hlt{$C^1$ Uses the Operator-Norm Metric}\\
Definition \ref{def:c1-map} makes $D\f$ itself the object whose continuity is measured.\\
For matrices $A,B\in\operatorname{Mat}(m,n,\R)$, $\abs{A_{ij}-B_{ij}}\leq\norm{A-B}_2\leq\norm{A-B}_F$. Moreover, $\norm{A-B}_F\leq\sqrt{mn}\max\limits_{i,j}\abs{A_{ij}-B_{ij}}$.\\
Thus operator-norm continuity and entrywise continuity are equivalent in these finite dimensions.
\end{perspective}

\noindent Figure \ref{fig:coordinate-polygon-partial-path} decomposes $\h$ into coordinate steps. This is the path used below to telescope from $\x$ to $\x+\h$.
\begin{illustrationfigure}[H]
  \centering
  \begin{tikzpicture}[
  scale=0.95,
  >=Stealth,
  every node/.style={font=\scriptsize},
  segment/.style={->, line width=0.85pt, blue!70!black},
  direct/.style={->, line width=0.55pt, black!40, dashed},
  point/.style={circle, fill=black, inner sep=1.15pt}
]
  \coordinate (p0) at (0,0);
  \coordinate (p1) at (1.55,0);
  \coordinate (p2) at (2.55,0.72);
  \coordinate (p3) at (2.55,2.02);

  \draw[direct] (p0) -- node[above left=2pt] {$\h$} (p3);
  \draw[segment] (p0) -- node[below=2pt] {$h_1\mathbf{e}_1$} (p1);
  \draw[segment] (p1) -- node[below right=1pt] {$h_2\mathbf{e}_2$} (p2);
  \draw[segment] (p2) -- node[right=2pt] {$h_3\mathbf{e}_3$} (p3);

  \node[point] at (p0) {};
  \node[point] at (p1) {};
  \node[point] at (p2) {};
  \node[point] at (p3) {};

  \node[left=3pt] at (p0) {$\x^{(0)}=\x$};
  \node[below=5pt] at (p1) {$\x^{(1)}$};
  \node[right=4pt] at (p2) {$\x^{(2)}$};
  \node[above=4pt] at (p3) {$\x^{(3)}=\x+\h$};
\end{tikzpicture}

  \caption{The solid coordinate steps sum to the dashed direct increment $\h$ from $\x$ to $\x+\h$.}
  \label{fig:coordinate-polygon-partial-path}
\end{illustrationfigure}

\begin{theorem} \label{thm:c1-iff-continuous-partials} \hlt{$C^1$ Means Continuous Partials}\\
Let $E\subseteq\R^n$ be open, $\f=(f_1,\ldots,f_m):E\to\R^m$. Then $\f$ is $C^1$ on $E$ iff $(\forall\,(i,j)\in J_m\times J_n)\,D_jf_i$ exists, continuous on $E$.
\begin{proof}
\begin{enumerate}[label=\roman*.]
\setlength{\itemsep}{0pt}
\item[$(\Rightarrow)$] Suppose $\f$ is $C^1$ on $E$. For $(\x,i,j)\in E\times J_m\times J_n$, $(D_jf_i)(\x)=[D\f(\x)]_{ij}$ by Theorem \ref{thm:total-derivative-matrix-partials}.\\
Since $\abs{A_{ij}-B_{ij}}\leq\norm{A-B}_2$, each $D_jf_i$ is continuous on $E$ by continuity of $D\f$.
\item[$(\Leftarrow)$] Fix $\x\in E$ and $i\in J_m$, and choose $\rho>0$ with $B_\rho(\x)\subseteq E$ by openness.\\
For $\norm{\h}_2<\rho$, set $\x^{(0)}\defeq\x$ and $\x^{(j)}\defeq\x+\sum\limits_{k=1}^{j}h_k\mathbf e_k$ for $j\in J_n$.\\
Each segment $[\x^{(j-1)},\x^{(j)}]$ lies in $B_\rho(\x)$ because every point on it is at distance at most $\norm{\h}_2$ from $\x$.\\
For $j\in J_n$ with $h_j\neq0$, let $\xi_j$ be a mean-value point on $[\x^{(j-1)},\x^{(j)}]$ by Corollary \ref{cor:mean-value-theorem}.\\
Then $f_i(\x^{(j)})-f_i(\x^{(j-1)})=h_j(D_jf_i)(\xi_j)$. For $j\in J_n$ with $h_j=0$, set $\xi_j\defeq\x^{(j)}$; the identity holds trivially.\\
Define $A_{i,\x}\h\defeq\sum\limits_{j=1}^{n}h_j(D_jf_i)(\x)$; then $A_{i,\x}\in\LL(\R^n,\R)$ by coordinatewise linearity.\\
The bound $\abs{f_i(\x+\h)-f_i(\x)-A_{i,\x}\h}\leq\sum\limits_{j=1}^{n}\abs{h_j}\abs{(D_jf_i)(\xi_j)-(D_jf_i)(\x)}$ holds by telescoping.\\
The right-hand side is at most $\sqrt n\norm{\h}_2\max\limits_{j\in J_n}\abs{(D_jf_i)(\xi_j)-(D_jf_i)(\x)}$ by Cauchy--Schwarz.\\
Here $\xi_j\to\x$ as $\h\to\ze$, and the maximum tends to $0$ by continuity and finiteness of $J_n$.\\
Hence displayed bound is $o(\norm{\h}_2)$. Thus every $f_i$ is differentiable at every $\x$, so $\f$ is differentiable on $E$ by Theorem \ref{thm:several-variable-componentwise-differentiability}(i).\\
Its derivative entries are $D_jf_i$ by Theorem \ref{thm:total-derivative-matrix-partials}. Continuity, $\norm{A-B}_2\leq\norm{A-B}_F\leq\sqrt{mn}\max\limits_{i,j}\abs{A_{ij}-B_{ij}}$ make $D\f$ continuous.\\
Hence $\f$ is $C^1$ on $E$ by Definition \ref{def:c1-map}.
\end{enumerate}
\hlt{Structural Commentary}: Telescoping builds the components; norm comparison assembles them.\\
\hlt{Key Insight}: Continuous partials control the finitely many coordinate remainders.
\end{proof}
\end{theorem}

Higher differentiation views $D\f$ as a map into the finite-dimensional space $\LL(\R^n,\R^m)$.\\
The second derivative is therefore a bilinear map, not merely an array of second partials.

\begin{definition} \label{def:higher-order-partials-and-c2} \hlt{Higher Derivatives and $C^r$}\\
Let $E\subseteq\R^n$ be open, $\f:E\to\R^m$. For $q\in\Z_{\geq1}$, $\LL^q(\R^n,\R^m)$ denote space of $q$-linear maps $(\R^n)^q\to\R^m$.\\
For $T\in\LL^q(\R^n,\R^m)$, set $\normop{T}\defeq\sup\limits_{\norm{\mathbf v_s}_2\leq1,\ s\in J_q} \norm{T(\mathbf v_1,\ldots,\mathbf v_q)}_2$.
\begin{enumerate}[label=(\roman*)]
\setlength{\itemsep}{0pt}
\item Set $D^1\f\defeq D\f$. For $q\geq2$, suppose $D^{q-1}\f$ is differentiable at $\x$.\\
The \hlt{$q$th derivative} is $D^q\f(\x)(\mathbf v_1,\ldots,\mathbf v_q)\defeq\bigl(D(D^{q-1}\f)(\x)\mathbf v_1\bigr)(\mathbf v_2,\ldots,\mathbf v_q)$.
\item For real-valued $f:E\to\R$, define the mixed second partial $D_{ij}f\defeq D_i(D_jf)$ wherever the right-hand side exists.
\item For $r\in\Z_{\geq1}$, the map $\f$ is \hlt{$C^r$} on $E$ iff $(\forall\,q\in J_r)\,D^q\f:E\to\LL^q(\R^n,\R^m)$ exists and is continuous.\\
The map $\f$ is \hlt{$C^\infty$} on $E$ iff $(\forall\,r\in\Z_{\geq1})\,\f$ is $C^r$ on $E$.
\end{enumerate}
\end{definition}

\begin{theorem} \label{thm:second-derivative-partials-and-continuity}
\hlt{Second Derivatives and Second Partials}\\
Let $E\subseteq\R^n$ be open, $\f=(f_1,\ldots,f_m):E\to\R^m$, $p\in E$. Set $\mathcal I\defeq J_n\times J_n\times J_m$.
\begin{enumerate}[label=(\roman*)]
\setlength{\itemsep}{0pt}
\item If $D^2\f(p)$ exists, then $(\forall\,(i,j,\ell)\in\mathcal I)\,D_{ij}f_\ell(p)$ exists. For these indices, $D^2 f_\ell(p)(\mathbf e_i,\mathbf e_j)=D_{ij}f_\ell(p)$.
\item Suppose $\f$ is differentiable on a neighborhood $U\subseteq E$ of $p$.\\
If $(\forall\,(i,j,\ell)\in\mathcal I)\,D_{ij}f_\ell$ exists on $U$ and is continuous at $p$, then $D^2\f(p)$ exists.
\end{enumerate}
\begin{proof}
\begin{enumerate}[label=(\roman*)]
\setlength{\itemsep}{0pt}
\item Identify $\LL(\R^n,\R^m)$ with $\R^{mn}$ by standard bases. Frobenius, operator norms equivalent because $\norm{M}_2\leq\norm{M}_F\leq\sqrt{mn}\norm{M}_2$.\\
The entries of $D\f$ are $D_jf_\ell$ by Theorem \ref{thm:total-derivative-matrix-partials}. Each entry is differentiable at $p$ by Theorem \ref{thm:several-variable-componentwise-differentiability}(i).\\
Its $i$th partial exists by Theorem \ref{thm:total-derivative-matrix-partials}(i). Consequently $Df_\ell$ is differentiable at $p$ by Theorem \ref{thm:several-variable-componentwise-differentiability}(i).\\
The identity $D^2 f_\ell(p)(\mathbf e_i,\mathbf e_j) =D_i(D_jf_\ell)(p)=D_{ij}f_\ell(p)$ holds by Definition \ref{def:higher-order-partials-and-c2}(i).
\item Set $F\defeq D\f:U\to\R^{mn}$ under the same identification, so $F_{\ell j}=D_jf_\ell$ by Theorem \ref{thm:total-derivative-matrix-partials}.\\
Define $A:\R^n\to\R^{mn}$ by $[A\h]_{\ell j}\defeq\sum\limits_{i=1}^{n}h_iD_{ij}f_\ell(p)$; then $A$ is linear. Choose $\rho>0$ with $B_\rho(p)\subseteq U$ by openness.\\
Define coordinate polygon $p^{(i)}\defeq p+\sum\limits_{q=1}^{i}h_q\mathbf e_q$. Set $S_i\defeq[p^{(i-1)},p^{(i)}]$; if $\norm{\h}_2<\rho$, then $S_i\subseteq B_\rho(p)$.\\
For $(i,j,\ell)\in\mathcal I$ with $h_i\neq0$, let $\eta_{i\ell j}$ be a mean-value point for $F_{\ell j}$ on $S_i$ by Corollary \ref{cor:mean-value-theorem}.\\
Then $F_{\ell j}(p^{(i)})-F_{\ell j}(p^{(i-1)})=h_iD_{ij}f_\ell(\eta_{i\ell j})$. For $(i,j,\ell)\in\mathcal I$ with $h_i=0$, set $\eta_{i\ell j}\defeq p^{(i)}$; the identity holds trivially.\\
For each $(\ell,j)$, $\abs{[F(p+\h)-F(p)-A\h]_{\ell j}}\leq\sum\limits_{i=1}^{n}\abs{h_i}\abs{D_{ij}f_\ell(\eta_{i\ell j})-D_{ij}f_\ell(p)}$ by telescoping.\\
This is at most $\sqrt n\norm{\h}_2\max\limits_{i\in J_n}\abs{D_{ij}f_\ell(\eta_{i\ell j})-D_{ij}f_\ell(p)}$ by Cauchy--Schwarz.\\
Since $\eta_{i\ell j}\to p$ as $\h\to\ze$, the right-hand side is $o(\norm{\h}_2)$ by continuity at $p$.\\
Therefore $\norm{F(p+\h)-F(p)-A\h}_2=o(\norm{\h}_2)$ by finiteness of $J_n\times J_m$.\\
This is also an operator-norm remainder by the preceding norm equivalence.\\
Thus $F=D\f$ is differentiable at $p$, and $D^2\f(p)$ exists by Definition \ref{def:higher-order-partials-and-c2}(i).
\end{enumerate}
\hlt{Structural Commentary}: Entrywise estimates combine into one derivative-matrix remainder.\\
\hlt{Key Insight}: Pointwise continuity suffices because only finitely many coordinates occur.
\end{proof}
\end{theorem}

\noindent Figure \ref{fig:signed-parallelogram-second-derivative} records the signed four-vertex difference used below.\\
Exchanging $\mathbf v$ and $\mathbf w$ preserves both the vertices and their signs.
\begin{illustrationfigure}[H]
  \centering
  \begin{tikzpicture}[
  scale=0.95,
  >=Stealth,
  every node/.style={font=\scriptsize},
  edge/.style={black!70, line width=0.7pt},
  stepvec/.style={->, blue!70!black, line width=0.75pt},
  point/.style={circle, fill=black, inner sep=1.2pt},
  sign/.style={font=\normalsize}
]
  \coordinate (p) at (0,0);
  \coordinate (v) at (4.1,0.35);
  \coordinate (w) at (1.15,2.25);
  \coordinate (vw) at (5.25,2.60);

  \fill[blue!4] (p) -- (v) -- (vw) -- (w) -- cycle;
  \draw[edge] (p) -- (v) -- (vw) -- (w) -- cycle;

  \node[point] at (p) {};
  \node[point] at (v) {};
  \node[point] at (w) {};
  \node[point] at (vw) {};

  \node[below left=2pt] at (p) {$p$};
  \node[below=4pt] at (v) {$p+t\mathbf{v}$};
  \node[above left=2pt] at (w) {$p+t\mathbf{w}$};
  \node[above right=2pt] at (vw) {$p+t\mathbf{v}+t\mathbf{w}$};

  \node[sign] at (0.62,0.38) {$+$};
  \node[sign] at (3.65,0.75) {$-$};
  \node[sign] at (1.42,1.90) {$-$};
  \node[sign] at (4.72,2.25) {$+$};

  \draw[stepvec] (0.45,-0.40) -- node[below=2pt] {$t\mathbf{v}$} (3.65,-0.13);
  \draw[stepvec] (-0.55,0.10) -- node[left=2pt] {$t\mathbf{w}$} (0.60,2.35);
\end{tikzpicture}

  \caption{The $+,-,-,+$ vertex pattern is invariant under exchanging $\mathbf v$ and $\mathbf w$.}
  \label{fig:signed-parallelogram-second-derivative}
\end{illustrationfigure}

\begin{theorem} \label{thm:second-derivative-symmetric} \hlt{Symmetry of the Second Derivative}\\
Let $E\subseteq\R^n$ be open, $\f:E\to\R^m$, $p\in E$. If $D^2\f(p)$ exists, then $(\forall\,\mathbf v,\mathbf w\in\R^n)\, D^2\f(p)(\mathbf v,\mathbf w)=D^2\f(p)(\mathbf w,\mathbf v)$.
\begin{proof}
\begin{enumerate}[label=(\roman*)]
\setlength{\itemsep}{0pt}
\item \hlt{Scalar Reduction}: Component projections commute with both derivatives by Theorem \ref{thm:several-variable-componentwise-differentiability}.\\
Suffices to prove symmetry for a scalar-valued $f$. Map $Df$ is defined on a neighborhood of $p$ by Definition \ref{def:higher-order-partials-and-c2}(i).
\item \hlt{Symmetric Increment}: Fix $\mathbf v,\mathbf w\in\R^n$. For $t$ sufficiently close to $0$, four arguments below lie in $E$ by openness.\\
Set $\Delta(t,\mathbf v,\mathbf w)\defeq f(p+t\mathbf v+t\mathbf w)-f(p+t\mathbf v)-f(p+t\mathbf w)+f(p)$.\\
This expression is unchanged when $\mathbf v$ and $\mathbf w$ are exchanged.
\item \hlt{Limit Formula}: Set $R(\mathbf z)\defeq Df(p+\mathbf z)-Df(p)-D^2f(p)(\mathbf z,\mathord\cdot)$. Then $\normop{R(\mathbf z)}=o(\norm{\mathbf z}_2)$ by Definition \ref{def:higher-order-partials-and-c2}(i).\\
For $t\neq0$, $\phi_t(s)\defeq f(p+t\mathbf v+s\mathbf w)-f(p+s\mathbf w)$. Let $\xi_t$ be mean-value point for $\phi_t$ between $0$ and $t$ by Corollary \ref{cor:mean-value-theorem}.\\
Then $\Delta(t,\mathbf v,\mathbf w)=t[Df(p+t\mathbf v+\xi_t\mathbf w)-Df(p+\xi_t\mathbf w)]\mathbf w$. Set $\mathbf z_0\defeq\xi_t\mathbf w$ and $\mathbf z_1\defeq t\mathbf v+\xi_t\mathbf w$.\\
Then $[Df(p+\mathbf z_1)-Df(p+\mathbf z_0)]\mathbf w=tD^2f(p)(\mathbf v,\mathbf w)+[R(\mathbf z_1)-R(\mathbf z_0)]\mathbf w$ by the definition of $R$.\\
Both remainder terms are $o(\abs t)$ because $\norm{\mathbf z_0}_2,\norm{\mathbf z_1}_2=O(\abs t)$. Therefore $D^2f(p)(\mathbf v,\mathbf w)=\lim\limits_{t\to0}\frac{\Delta(t,\mathbf v,\mathbf w)}{t^2}$.
\item \hlt{Symmetry}: Same formula with $\mathbf v$, $\mathbf w$ exchanged has same numerator by part (ii). Hence $D^2f(p)(\mathbf v,\mathbf w)=D^2f(p)(\mathbf w,\mathbf v)$.
\end{enumerate}
\hlt{Structural Commentary}: A symmetric four-point increment encodes either bilinear order.\\
\hlt{Key Insight}: One geometric numerator computes both orders of directional differentiation.
\end{proof}
\end{theorem}

\begin{corollary} \label{cor:mixed-partials-from-second-derivative}
\hlt{Mixed Partials from $D^2$}\\
Let $E\subseteq\R^n$ be open, $\f=(f_1,\ldots,f_m):E\to\R^m$, $p\in E$. If $D^2\f(p)$ exists, then $(\forall\,(i,j,\ell)\in J_n\times J_n\times J_m)\, D_{ij}f_\ell(p)=D_{ji}f_\ell(p)$.
\begin{proof}
For $(i,j,\ell)\in J_n\times J_n\times J_m$, $D_{ij}f_\ell(p)=D^2f_\ell(p)(\mathbf e_i,\mathbf e_j)$ by Theorem \ref{thm:second-derivative-partials-and-continuity}(i).\\
Likewise, $D_{ji}f_\ell(p)=D^2f_\ell(p)(\mathbf e_j,\mathbf e_i)$ by the same part. The two right-hand sides are equal by Theorem \ref{thm:second-derivative-symmetric}.\\
\hlt{Structural Commentary}: Second partials are standard-basis values of one bilinear derivative.\\
\hlt{Key Insight}: Symmetry of $D^2f_\ell(p)$ identifies the two mixed entries.
\end{proof}
\end{corollary}

\begin{corollary} \label{cor:higher-derivatives-symmetric} \hlt{Symmetry of Higher Derivatives}\\
Let $E\subseteq\R^n$ be open, $\f:E\to\R^m$, $p\in E$, $r\in\Z_{\geq2}$, $\mathbf v=(\mathbf v_1,\ldots,\mathbf v_r)\in(\R^n)^r$.\\
For $\sigma\in\mathfrak S_r$, set $\mathbf v^\sigma\defeq(\mathbf v_{\sigma(1)},\ldots,\mathbf v_{\sigma(r)})$. If $D^r\f(p)$ exists, then $(\forall\,\sigma\in\mathfrak S_r)\,D^r\f(p)(\mathbf v^\sigma)=D^r\f(p)(\mathbf v)$.
\begin{proof}
\begin{enumerate}[label=(\roman*)]
\setlength{\itemsep}{0pt}
\item \hlt{Base Case}: The case $r=2$ holds by Theorem \ref{thm:second-derivative-symmetric}.
\item \hlt{Trailing Slots}: Let $r>2$, assume claim for order $r-1$. Since $D^r\f(p)$ exists, map $D^{r-1}\f$ is defined on neighborhood of $p$.\\
Near $p$, each value of $D^{r-1}\f$ is symmetric by the induction hypothesis.\\
The symmetric $(r-1)$-linear maps form a closed finite-dimensional subspace.\\
Thus $D(D^{r-1}\f)(p)\mathbf v_1$ is symmetric in slots $2,\ldots,r$ by closure under limits.
\item \hlt{First Two Slots}: Identify the finite-dimensional codomain of $D^{r-2}\f$ with a Euclidean space.\\
Its second derivative at $p$ exists by the existence of $D^r\f(p)$ and is symmetric by Theorem \ref{thm:second-derivative-symmetric}.\\
After evaluation at $\mathbf v_3,\ldots,\mathbf v_r$, $D^r\f(p)$ is symmetric in slots $1$ and $2$ by the preceding identity.
\item \hlt{Permutation Step}: Invariance under every adjacent transposition holds by parts (ii)--(iii).\\
Since adjacent transpositions generate $\mathfrak S_r$, every permutation fixes $D^r\f(p)$.
\end{enumerate}
\hlt{Structural Commentary}: Induction controls the tail; $D^2$ symmetry swaps the first pair.\\
\hlt{Key Insight}: Adjacent transpositions generate every permutation of the derivative slots.
\end{proof}
\end{corollary}

The next argument uses coordinate rectangles to compare the two iterated partial derivatives.\\
Its local continuity hypothesis is weaker than a global $C^2$ assumption.

\begin{theorem} \label{thm:rectangle-mean-value-mixed-partial} \hlt{Rectangle Mean Value Formula}\\
Let $E\subseteq\R^2$ be open, $f:E\to\R$, and suppose $D_1f$ and $D_{21}f=D_2(D_1f)$ exist on $E$.\\
Let $h,k\in\R\setminus\{0\}$, $Q\subseteq E$ be closed coordinate rectangle with opposite vertices $(a,b)$ and $(a+h,b+k)$.\\
Set $\Delta(f,Q)\defeq f(a+h,b+k)-f(a+h,b)-f(a,b+k)+f(a,b)$. Then $(\exists\,(x,y)\in\operatorname{int}Q)\, \Delta(f,Q)=hk(D_{21}f)(x,y)$.
\begin{proof}
Let $I$ and $J$ be the open intervals with endpoints $a,a+h$ and $b,b+k$, respectively. Set $u(t)\defeq f(t,b+k)-f(t,b)$.\\
The map $u$ is differentiable on $I$ and continuous on its closure by existence of $D_1f$.\\
Let $x\in I$ be a mean-value point for $u$ by Corollary \ref{cor:mean-value-theorem}. Then $\Delta(f,Q)=h[(D_1f)(x,b+k)-(D_1f)(x,b)]$.\\
The map $s\mapsto(D_1f)(x,s)$ is differentiable on $J$ and continuous on its closure by existence of $D_{21}f$.\\
Let $y\in J$ be a mean-value point for $s\mapsto(D_1f)(x,s)$ by the same corollary.\\
Then $(D_1f)(x,b+k)-(D_1f)(x,b)=k(D_{21}f)(x,y)$. Then $(x,y)\in I\times J=\operatorname{int}Q$ and $\Delta(f,Q)=hk(D_{21}f)(x,y)$.\\
\hlt{Structural Commentary}: Two mean-value steps factor the rectangle increment.\\
\hlt{Key Insight}: The four-corner difference samples the vertical change of $D_1f$.
\end{proof}
\end{theorem}

\begin{theorem} \label{thm:equality-of-mixed-partials} \hlt{Equality of Mixed Partials}\\
Let $E\subseteq\R^2$ be open, $f:E\to\R$, $(a,b)\in E$. Suppose $D_1f$, $D_2f$, and $D_{21}f$ exist on $E$, $D_{21}f$ is continuous at $(a,b)$. Then:
\begin{enumerate}[label=(\roman*)]
\setlength{\itemsep}{0pt}
\item The mixed partial $(D_{12}f)(a,b)$ exists.
\item $(D_{12}f)(a,b)=(D_{21}f)(a,b)$.
\end{enumerate}
\begin{proof}
Set $A\defeq(D_{21}f)(a,b)$.
\begin{enumerate}[label=(\roman*)]
\setlength{\itemsep}{0pt}
\item Fix $\varepsilon>0$, $\delta>0$ with $B_\delta(a,b)\subseteq E$ by openness. After shrinking $\delta$, $(\forall\,\mathbf z\in B_\delta(a,b))\,\abs{(D_{21}f)(\mathbf z)-A}<\frac{\varepsilon}{2}$ by continuity at $(a,b)$.\\
For $h,k\neq0$ with $\sqrt{h^2+k^2}<\delta$, let $Q_{h,k}$ be coordinate rectangle. Its opposite vertices are $(a,b)$ and $(a+h,b+k)$.\\
Then $Q_{h,k}\subseteq B_\delta(a,b)\subseteq E$. Let $\mathbf z_{h,k}\in\operatorname{int}Q_{h,k}$ be a rectangle mean-value point by Theorem \ref{thm:rectangle-mean-value-mixed-partial}.\\
Then $\frac{\Delta(f,Q_{h,k})}{hk}=(D_{21}f)(\mathbf z_{h,k})$. Hence $\abs{\frac{\Delta(f,Q_{h,k})}{hk}-A}<\frac{\varepsilon}{2}$. Fix $h$ with $0<\abs h<\delta$ and let $k\to0$.\\
Since $D_2f$ exists at $(a+h,b)$ and $(a,b)$, $\frac{\Delta(f,Q_{h,k})}{hk}\to \frac{(D_2f)(a+h,b)-(D_2f)(a,b)}{h}$; last quotient differs from $A$ by at most $\frac{\varepsilon}{2}<\varepsilon$.\\
This holds for every $h$ with $0<\abs h<\delta$, so the quotient tends to $A$ as $h\to0$. Thus $(D_{12}f)(a,b)$ exists.
\item The limit in part (i) equals $A=(D_{21}f)(a,b)$ by the definition of $A$.
\end{enumerate}
\hlt{Structural Commentary}: Continuity traps each shrinking rectangle quotient near one value.\\
\hlt{Key Insight}: The vertical limit recovers $D_2f$; the horizontal limit produces $D_{12}f$.
\end{proof}
\end{theorem}

\begin{corollary} \label{cor:c2-mixed-partials-commute} \hlt{$C^2$ Mixed Partials Commute}\\
Let $E\subseteq\R^2$ be open, $f:E\to\R$ be $C^2$. Then $(\forall\,(i,j)\in J_2\times J_2)\,D_{ij}f=D_{ji}f$ on $E$.
\begin{proof}
Fix $p\in E$, $(i,j)\in J_2\times J_2$. Derivative $D^2f(p)$ exists by Definition \ref{def:higher-order-partials-and-c2}(iii); $D_{ij}f(p)=D_{ji}f(p)$ by Corollary \ref{cor:mixed-partials-from-second-derivative}.\\
\hlt{Structural Commentary}: The $C^2$ hypothesis ensures second differentiability at every point.\\
\hlt{Key Insight}: Pointwise symmetry of $D^2f$ makes the coordinate partials commute globally.
\end{proof}
\end{corollary}

\begin{perspective} \label{per:mixed-partials-need-hypotheses}
\hlt{Mixed Partials Need Hypotheses}\\
The quantities $D_{12}f$ and $D_{21}f$ arise from different iterated limits.\\
Existence of $D^2f$ guarantees equality by Corollary \ref{cor:mixed-partials-from-second-derivative}. Continuity of one mixed partial guarantees it by Theorem \ref{thm:equality-of-mixed-partials}.
\end{perspective}
